\documentclass[border=5mm]{standalone}
% ==============================================================================
% SETUP.TEX - CONFIGURACIÓN GLOBAL DEL PROYECTO
% ==============================================================================
% Este archivo centraliza todos los paquetes y estilos.
% Debe incluirse en el preámbulo del libro principal y de los archivos standalone.
% \PassOptionsToPackage{gray}{xcolor}
% --- 1. CODIFICACIÓN E IDIOMA ---
% fix-cm habilita las fuentes Computer Modern en tamaños arbitrarios (por debajo
% de 5 pt, entre otros). Debe cargarse antes que fontenc.
\usepackage{fix-cm}
\usepackage[utf8]{inputenc}
\usepackage[T1]{fontenc}
\usepackage[spanish, es-tabla]{babel}  
\usepackage{csquotes} % Recomendado para babel y citas

\usepackage{import}
% mode=image: Usa el PDF pre-compilado, nunca intenta recompilar.
% Debes compilar cada figura manualmente antes de compilar el libro.
\usepackage[mode=image, subpreambles=false]{standalone}

% --- 2. MATEMÁTICAS Y CIENCIAS ---
\usepackage{amsmath, amssymb, amsfonts, mathtools}
\usepackage{enumitem}

% LaTeX trae \DeclareMathSizes{5}{5}{5}{5}, así que a 5 pt (\tiny) los subíndices
% salen del mismo tamaño que la base: en las etiquetas de figuras el M de
% $\omega_M$ queda desproporcionado. Se restablece la proporción habitual.
\DeclareMathSizes{5}{5}{3.7}{3}

% --- 3. DISEÑO Y GEOMETRÍA --- 
% Unificamos los márgenes para todo el documento
\usepackage[left=2.5cm,right=2.5cm,top=3cm,bottom=3cm]{geometry}
\makeatletter
\ifstandalone % Evita que geometry dañe el recorte de las figuras bajándolas 3cm
  \@ifclasswith{standalone}{crop=false}{
    % Es un capítulo/sección: Dejamos que geometry actúe.
  }{
    % Es una figura: Neutralizamos geometry para que no las recorte mal.
    \geometry{pass}
  }
\fi
\makeatother
\usepackage{parskip} % Espaciado entre párrafos sin sangría
\usepackage{float}   % Para usar [H] en figura

% Ajustes de flotantes para evitar espacios verticales exagerados
\renewcommand{\topfraction}{0.95}      % Permite que las figuras ocupen hasta el 95% superior
\renewcommand{\bottomfraction}{0.95}   % Permite que las figuras ocupen hasta el 95% inferior
\renewcommand{\textfraction}{0.05}     % Permite hojas con tan solo un 5% de texto
\renewcommand{\floatpagefraction}{0.8} % Requiere que una página de solo figuras esté 80% llena
\setcounter{topnumber}{4}
\setcounter{bottomnumber}{4}
\setcounter{totalnumber}{10}

% --- 4. GRÁFICOS Y TABLAS ---
\usepackage{graphicx}
\usepackage{booktabs} % Tablas profesionales
\usepackage{caption}  % Control de captions y \captionof
\usepackage{tikz}
\usetikzlibrary{arrows.meta, calc, angles, quotes, babel, backgrounds}
\usepackage{pgfplots}
\usepgfplotslibrary{groupplots}
\usepgfplotslibrary{external}
\usepgfplotslibrary{patchplots}
\pgfplotsset{compat=1.18}
\usepackage[american, straightvoltages]{circuitikz}

% --- 5. COLORES Y CAJAS ---

\usepackage[table]{xcolor}
\usepackage{tcolorbox}

% Definición de colores corporativos/estilo
\definecolor{utnblue}{RGB}{0, 51, 102}
\definecolor{codegreen}{rgb}{0,0.6,0}
\definecolor{codegray}{rgb}{0.5,0.5,0.5}
\definecolor{codepurple}{rgb}{0.58,0,0.82}
\definecolor{backcolour}{rgb}{0.96,0.96,0.96}

% --- 6. ENTORNOS PERSONALIZADOS (CAJAS) ---

% Caja "Resultado" (Azul Corporativo) — Definiciones, fórmulas clave, resultados importantes.
% Uso: \begin{resultado}[Fórmula de Euler] ... \end{resultado}
\newtcolorbox{resultado}[1][]{
    colback=utnblue!5!white,
    colframe=utnblue,
    title=\textbf{#1},
    fonttitle=\bfseries,
    sharp corners,
    boxrule=0.5mm
}

% Caja "Nota" (Gris) — Advertencias, aplicaciones, contexto secundario.
% Uso: \begin{nota}[Cuidado con la calculadora] ... \end{nota}
% Uso: \begin{nota} ... \end{nota}  → título por defecto "Nota"
\newtcolorbox{nota}[1][Nota]{
    colback=codegray!5!white,
    colframe=codegray,
    title=\textbf{#1},
    fonttitle=\bfseries,
    sharp corners,
    boxrule=0.5mm
}

% Caja inline para resaltar un valor y enlazarlo con su otra aparición en una tabla
% o en una cuenta: mismo color, mismo valor.
% Uso: \marcavalor{$4$}  o  \marcavalor[codegreen]{$4$}
\newtcbox{\marcavalor}[1][utnblue]{
    on line,
    boxsep=1pt, left=2pt, right=2pt, top=1pt, bottom=1pt,
    colback=#1!10!white,
    colframe=#1,
    boxrule=0.4pt,
    arc=2pt
}

% --- 7. CÓDIGO FUENTE (LISTINGS) ---
\usepackage{listings}

\lstdefinestyle{miestilo}{
    backgroundcolor=\color{backcolour},   
    commentstyle=\color{codegreen},
    keywordstyle=\color{magenta},
    numberstyle=\tiny\color{codegray},
    stringstyle=\color{codepurple},
    basicstyle=\ttfamily\footnotesize,
    breakatwhitespace=false,         
    breaklines=true,                 
    captionpos=b,                    
    keepspaces=true,                 
    numbers=left,                    
    numbersep=5pt,                  
    showspaces=false,                
    showstringspaces=false,
    showtabs=false,                  
    tabsize=2,
    frame=single,                    
    rulecolor=\color{codegray}
}

\lstset{style=miestilo}
% Soporte para tildes y ñ en bloques de código
\lstset{literate={ñ}{{\~n}}1 {á}{{\'a}}1 {é}{{\'e}}1 {í}{{\'i}}1 {ó}{{\'o}}1 {ú}{{\'u}}1}

% Operadores matemáticos personalizados

% Vector en sentido discreto (lista finita de coordenadas): negrita + flecha.
% Uso: \vect{v}, \vect{e}_k, \vect{0}.
\newcommand{\vect}[1]{\vec{\mathbf{#1}}}

% Fecha de versión y comando para capítulos standalone
\providecommand{\verdate}{\the\year.\ifnum\month<10 0\fi\the\month.\ifnum\day<10 0\fi\the\day}
\newcommand{\chapterversion}{%
    \vspace{-1.5cm}\begin{flushright}\small\textit{\color{codegray}Versión~\verdate}\end{flushright}\vspace{0.5cm}%
}

% --- 8. HIPERVÍNCULOS (DEBE SER EL ÚLTIMO PAQUETE) ---
\usepackage[colorlinks=true, linkcolor=utnblue, urlcolor=utnblue, citecolor=utnblue]{hyperref}

% --- 9. REFERENCIAS CRUZADAS ENTRE CAPÍTULOS STANDALONE ---
% Cuando se compila un capítulo standalone, xr-hyper lee los .aux de los
% otros capítulos (si existen) para resolver \ref{} a labels externos.
% En la compilación del libro completo este bloque no se activa.
\ifstandalone
  \usepackage{xr-hyper}
  \makeatletter
  \newcommand{\xrchap}[1]{%
    \edef\xrc@self{\detokenize{Capitulo#1}}%
    \edef\xrc@job{\jobname}%
    \ifx\xrc@self\xrc@job\else
      \IfFileExists{../#1capitulo/Capitulo#1.aux}%
        {\externaldocument{../#1capitulo/Capitulo#1}}{}%
    \fi
  }
  \makeatother
  \xrchap{01}\xrchap{02}\xrchap{03}\xrchap{04}
  \xrchap{05}\xrchap{06}\xrchap{07}\xrchap{08}
  \xrchap{09}\xrchap{10}\xrchap{11}\xrchap{12}
  \xrchap{13}
\fi
% \PassOptionsToPackage{gray}{xcolor}

\begin{document}
    \begin{tikzpicture}[>=Latex, font=\scriptsize,
        mk/.style={fill=utnblue},
        vec/.style={utnblue, semithick, -{Latex[length=4pt]}},
        arco/.style={red!80!black, semithick, -{Latex[length=3pt]}},
    ]
        \def\rc{1.15}    % radio del círculo
        \def\rl{1.42}    % radio de las etiquetas de muestra
        \def\ra{0.48}    % radio del arco que mide lambda
        \def\dcol{5.80}  % paso horizontal entre paneles
        \def\dfil{4.50}  % paso vertical entre filas

        \newcommand{\circulo}{\draw[gray!55, thin] (0,0) circle (\rc);}
        % muestras sobre el círculo, numeradas con el índice n
        % #1 = radio del marcador, #2 = lista angulo/etiqueta
        \newcommand{\puntos}[2]{%
            \foreach \ang/\lbl in {#2} {
                \fill[mk] (\ang:\rc) circle (#1);
                \node[utnblue, font=\tiny, inner sep=1pt] at (\ang:\rl) {\lbl};
            }%
        }
        % marcador de panel y valor de lambda, juntos arriba y centrados
        \newcommand{\panel}[2]{%
            \node[anchor=south] at (0,1.66) {#1\quad$#2$};
        }

        % ============================== FILA SUPERIOR ==============================

        % ---------- (a) k = 0 ----------
        \begin{scope}[shift={(0,0)}]
            \panel{(a)}{\lambda_0 = 0}
            \circulo
            \puntos{1.8pt}{0/0}
            \draw[vec] (0,0) -- (0:\rc);
        \end{scope}

        % ---------- (b) k = 1 ----------
        \begin{scope}[shift={(\dcol,0)}]
            \panel{(b)}{\lambda_1 = \pi/4}
            \circulo
            \puntos{1.6pt}{0/0, 45/1, 90/2, 135/3, 180/4, 225/5, 270/6, 315/7}
            \draw[vec] (0,0) -- (45:\rc);
            \draw[arco] (0:\ra) arc (0:45:\ra);
        \end{scope}

        % ---------- (c) k = 2 ----------
        \begin{scope}[shift={(2*\dcol,0)}]
            \panel{(c)}{\lambda_2 = \pi/2}
            \circulo
            \puntos{1.6pt}{0/0, 90/1, 180/2, 270/3}
            \draw[vec] (0,0) -- (90:\rc);
            \draw[arco] (0:\ra) arc (0:90:\ra);
        \end{scope}

        % ---------- (d) k = 3 ----------
        \begin{scope}[shift={(3*\dcol,0)}]
            \panel{(d)}{\lambda_3 = 3\pi/4}
            \circulo
            \puntos{1.6pt}{0/0, 135/1, 270/2, 45/3, 180/4, 315/5, 90/6, 225/7}
            \draw[vec] (0,0) -- (135:\rc);
            \draw[arco] (0:\ra) arc (0:135:\ra);
        \end{scope}

        % ============================== FILA INFERIOR ==============================

        % ---------- (e) k = 4 ----------
        \begin{scope}[shift={(0,-\dfil)}]
            \panel{(e)}{\lambda_4 = \pi}
            \circulo
            \puntos{1.8pt}{0/0, 180/1}
            \draw[vec] (0,0) -- (180:\rc);
            \draw[arco] (0:\ra) arc (0:180:\ra);
        \end{scope}

        % ---------- (f) k = 5 ----------
        \begin{scope}[shift={(\dcol,-\dfil)}]
            \panel{(f)}{\lambda_5 = 5\pi/4}
            \circulo
            \puntos{1.6pt}{0/0, 225/1, 90/2, 315/3, 180/4, 45/5, 270/6, 135/7}
            \draw[vec] (0,0) -- (225:\rc);
            \draw[arco] (0:\ra) arc (0:225:\ra);
        \end{scope}

        % ---------- (g) k = 6 ----------
        \begin{scope}[shift={(2*\dcol,-\dfil)}]
            \panel{(g)}{\lambda_6 = 3\pi/2}
            \circulo
            \puntos{1.6pt}{0/0, 270/1, 180/2, 90/3}
            \draw[vec] (0,0) -- (270:\rc);
            \draw[arco] (0:\ra) arc (0:270:\ra);
        \end{scope}

        % ---------- (h) k = 7 ----------
        \begin{scope}[shift={(3*\dcol,-\dfil)}]
            \panel{(h)}{\lambda_7 = 7\pi/4}
            \circulo
            \puntos{1.6pt}{0/0, 315/1, 270/2, 225/3, 180/4, 135/5, 90/6, 45/7}
            \draw[vec] (0,0) -- (315:\rc);
            \draw[arco] (0:\ra) arc (0:315:\ra);
        \end{scope}

    \end{tikzpicture}
\end{document}
